\section{Stationarity directly from the sequential equilibrium}\label{sec:stationarity}

The usual route to a stationary equilibrium passes through a recursive state space, often enlarged
and randomised \citep{dgmm1994}. A sequence-state equilibrium offers a direct one: it is a function
of the infinite history of shocks, and the history evolves exogenously, so if the history is drawn
from a law left invariant by the shock process, the aggregate state $(z,\mu^*)$ has the same law
every period. The equilibrium is a factor of the stationary shock process.

\begin{theorem}\label{thm:factor}
The shock-history process has an invariant ergodic law, unique when the aggregate chain is bounded
below. For every sequence-state equilibrium measurable in the history --- with no uniqueness ---
tomorrow's aggregate state has today's law. Under uniqueness of sequential equilibrium that law is
an invariant distribution of the equilibrium Markov process on $(z,\mu)$, ergodic when the shock law
is.
\end{theorem}
\formal{exists\_invariant\_ergodic, stationary\_unique\_of\_minor, stateLaw\_tomorrow,
isInvariant\_stateLaw}

\begin{theorem}\label{thm:stationary}
For CRRA utility, Cobb--Douglas prices with $0\le\delta<1$ and $(\beta(1-\delta))^{1/\gamma}<1-\delta$,
both shock chains bounded below and a small enough capital share: the sequence-state equilibrium is
unique and continuous in the history; the law of its aggregate state under the stationary shock law
is an invariant and ergodic distribution of the equilibrium Markov process; and every invariant
distribution of that process equals it.
\end{theorem}
\formal{continuous\_seqState, minMPC\_smallShare\_unique\_stationary}

\noindent
Continuity comes from fading memory (Theorem~\ref{thm:matching}) and Stone--Weierstrass density of
asset-Lipschitz test functions. Uniqueness of the invariant distribution is a coupling: the $T$-step
expectation of a test function from any state is within $C(r^*)^T$ of a function of today's shock
alone, invariant distributions have the same shock marginal, and the test functions separate
measures on the Polish aggregate state space. The theorem proves uniqueness of the invariant
distribution, not convergence in law from every starting point.
